\documentclass{article}

% Symbols
%\usepackage{recycle}
\usepackage{amsfonts, amsthm}
\usepackage{upgreek}
\usepackage{physics}
\usepackage{cancel}
\usepackage{amssymb, latexsym, amsmath}

% Proof
\renewcommand*{\proofname}{\textbf{Demostraci\'on:}}

% Graphics
\usepackage{graphicx}
\usepackage{pgf}

% Color a letras.
%\usepackage[usenames,dvipsnames,svgnames,table]{xcolor}

% Tikz
\usepackage{tkz-graph}
\usepackage{tikz}
\usetikzlibrary{arrows,automata}
\usepackage{tikz}
\usetikzlibrary{arrows,automata}
% Def. Dr. César.
\usetikzlibrary{shapes,calc}
\tikzstyle{edge}=[shorten <=2pt, shorten >=2pt, >=stealth, line width=1.1pt]
\tikzstyle{blueE}=[shorten <=2pt, shorten >=2pt, >=stealth, line width=1.5pt, blue]
\tikzstyle{blackV}=[circle, fill=black, minimum size=6pt, inner sep=0pt, outer sep=0pt]
\tikzstyle{blueV}=[circle, fill=blue, draw, minimum size=6pt, line width=0.75pt, inner sep=0pt, outer sep=0pt]
\tikzstyle{redV}=[circle, fill=red, draw, minimum size=6pt, line width=0.75pt, inner sep=0pt, outer sep=0pt]
\tikzstyle{redSV}=[semicircle, fill=red, minimum size=3pt, inner sep=0pt, outer sep=0pt, rotate=225]
\tikzstyle{blueSV}=[semicircle, fill=blue, minimum size=3pt, inner sep=0pt, outer sep=0pt, rotate=225]
\tikzstyle{blackSV}=[semicircle, fill=black, minimum size=3pt, inner sep=0pt, outer sep=0pt, rotate=225]
\tikzstyle{vertex}=[circle, draw, minimum size=6pt, line width=0.75pt, inner sep=0pt, outer sep=0pt]

% Nueva Def.
\newcommand{\kaspressknodel}[2] {
  \pgfmathtruncatemacro{\mymod}{#1+1}
  \foreach \X in {0,...,#1} {
    \node[bullet,label={right:$\boldsymbol{(2,\X)}$}] (v2\X) at (4,{1.2*(\X-(#1+0)/2)}){};
    \node[bullet,label={left:$\boldsymbol{(1,\X)}$}] (v1\X) at (1.5,{1.2*(\X-(#1+0)/2)}){};
  }
  \foreach \X [evaluate=\X as \NextX using {int(mod(\X+3,\mymod))},
               evaluate=\X as \AnotherX using {int(mod(\X+1,\mymod))},] in {0,...,#1} {
    \draw[thick] (v1\X) -- (v2\NextX);
    \draw[dashed] (v1\X) -- (v2\AnotherX);
    \draw (v1\X) -- (v2\X);
  }
}

% Margins
\addtolength{\voffset}{-1.5cm}
\addtolength{\hoffset}{-1.5cm}
\addtolength{\textwidth}{3cm}
\addtolength{\textheight}{3cm}

%Header-Footer
\usepackage{fancyhdr}
\renewcommand{\headrulewidth}{1pt}

\newcommand{\set}[1]{%
  \left\{ #1 \right\}%
}

%\pagenumbering{gobble} -- Este comando
%                       -- quita el número de página.
\footskip = 50pt
\renewcommand{\headrulewidth}{1pt}

\pagestyle{fancyplain}

\begin{document}
\title{UNIVERSIDAD AUT\'ONOMA DE M\'EXICO\\ Facultad de Ciencias}
\author{Autores:
  \\ Fernanda Villaf\'an Flores
  \\ Fernando Alvarado Palacios
  \\ Adri\'an Aguilera Moreno}
\date{}
\maketitle
\begin{center}
  \includegraphics[scale=0.20]{../Imagen/Portada.jpg}\\[0.4cm]
  \Large
  \bf{Gr\'aficas y Juegos}
  \normalsize
\end{center}
\newpage
\fancyhead[r]{ Gr\'aficas y Juegos 2022-1}
%%%%%%%%%%%%%%%%%%%%%%%%%%%%%%%%%%%%%%%%%%%%%%%%%%%%%
\section*{\LARGE{Tarea 4}}

\begin{enumerate}
  %%%%%%%%%%%%%%%%%%%%%%%%%%%%%%%%%%%% Ejercicio 1. %%%%%%%%%%%%%%%%%%%%%%%%%%%%%
  \item Sea $G$ una gr\'afica no trivial.   Demuestre
    que $G$ es una trayectoria si y s\'olo si $G$ es
    un \'arbol con exactamente dos v\'ertices de
    grado $1$.

    \begin{proof}
      $\Longrightarrow$) Sea $G$ una trayectoria.

      Entonces para cualesquiera dos vértices $u, v$ que pertenecen a $G$, $u, v$ son
      adyacentes, es decir, $u, v$ son consecutivos.

      Sea $v_t$ el vértice inicial y $v_q$ el vértice final de la trayectoria.

      Así, $v_t$ y $v_q$ serán de grado $1$ ya que no existe un vértice antes o
      uno después.

      Por lo tanto, existen únicamente $2$ vértices en $G$ con grado $1$ ya que
      los demás son internos a la trayectoria.

      $\Longleftarrow$) Sea  $V_G$ el conjunto de vértices de $G$.

      Sean $v_r$, $v_s$ cualesquiera vértices pertenecientes a $V_G$ tales que
      $v_r$ y $v_s$ tienen grado $1$.

      Como $G$ es un árbol, entonces todos los vértices de $G$ están relacionados
      y tampoco existe ningún ciclo en $G$. Así, $G$ admite un orden lineal
      empezando la trayectoria (s.p.g.) por $v_r$ y finalizando en $v_s$.

      Por tanto, si $u,v$ son adyacentes entonces $u,v$ son consecutivos y
      concluimos que $G$ es una trayectoria. \\
    \end{proof}

    %%%%%%%%%%%%%%%%%%%%%%%%%%%%%%% Ejercicio 2. %%%%%%%%%%%%%%%%%%%%%%%%%%%%%%%
    \item \begin{enumerate}
      %~~~~~~~~~~~~~~~~~~~~~~~~~~~~~~~~~~~~~ 2 (a).
      \item Demuestre que cada \'arbol con grado m\'aximo
        $\Delta > 1$ tiene al menos $\Delta$ hojas.

        \renewcommand\qedsymbol{QED}
        \begin{proof}
          Sea $G$ un \'arbol y sea $x \in V_G$ tal que $d(x) = \Delta$ (notar que $x$
          no necesariamente es el \'unico que cumple tener grado igual a $\Delta$, por
          lo que se tomar\'a alguno que cumpla esto), entonces $x$ tiene exactamente
          $\Delta$ vecinos.

          Por la caracterizaci\'on de \'arbol, sabemos que $G$ es ac\'iclico
          y por tanto, los caminos que parten desde $x$ (tienen a $x$
          como v\'ertice inicial) hacia algunos de sus $\Delta$ vecinos no tienen
          v\'ertices en com\'un que sean distintos de $x$. En caso contrario, habr\'ia $2$
          trayectorias que tienen inicio en $x$ a las que llamaremos $w_1$ y $w_2$, y
          adem\'as tienen en com\'un al menos un v\'ertice $v$ y por tanto, $w_1 w_2$ es
          un ciclo!! que est\'a contenido en $G$.

          Luego, como $x$ tiene $\Delta$ vecinos, entonces podemos tomar al menos
          $\Delta$ trayectorias distintas entre ellas tales que terminen en alg\'un
          v\'ertice $u_i$ ($1 \leq i \leq \Delta$).

          Así, cada $u_i$ es una hoja de $G$ y como hemos encontrado $\Delta$ hojas
          podemos concluir que el \'arbol $G$ con grado m\'aximo $\Delta$ tiene al menos
          $\Delta$ hojas. \\
          \\
        \end{proof}
  \newpage
      %~~~~~~~~~~~~~~~~~~~~~~~~~~~~~~~~~~~ 2 (b).
      \item Construya, para cada elecci\'on de $n$ y $\Delta$,
        con $2\le \Delta < n$, un \'arbol de orden $n$ con
        exactamente $\Delta$ hojas.

          \textbf{\textit{Soluci\'on:}}

          Sea $G$ un \'arbol.

          Usando el resultado anterior, garantizamos que $G$ tiene alg\'un
          $x \in V_G : d(x) = \Delta$, entonces $G$ tiene al menos $\Delta$ hojas.

          Luego, part\'icularmente en los \'arboles de
          exactamente $\Delta$ hojas, cada camino $W_i$ que tenga como v\'ertice
          inicial a $x$ no tendr\'a bifurcaciones. Esto es, $x W_i v_i$ es la
          \'unica manera de llegar de $x$ a $v_i$ y adem\'as $W_i$ es una trayectoria
          (y por definici\'on de trayectoria, no tendr\'a bifurcaciones).

          Como $x$ es v\'ertice inicial de exactamente $\Delta$ trayectorias $W_i$,
          entonces hay exactamente $\Delta$ hojas en $G$.

          En resumen, $G$ tendr\'a exactamente $\Delta$ v\'ertices de grado $1$,
          a continuaci\'on se muestra un ejemplo que trata de ser lo m\'as general
          posible:

      %%%%%%%%%%%%%%%%%%%%%%%%%%%%%%%%%%%%%%%%%%%%%%%%%%%%%%%%%%%%%%%%%%%%%%%%%%
        \begin{figure}[ht!]
          \centering
          \begin{tikzpicture}
            %%%%%  Componente izquierda  %%%%%
            %\foreach \i in {0,1}
            %\node(\i) [vertex] at (\i*1, -2){};
            \node(0) [blueV, label=270:$v_1$] at (0, -2){};
            \node(1) [vertex] at (1, -2){};
            \draw[edge] (0) to (1);
            \node (L) at (-0.5,-1.5){$G$};

            %% Suspensivos.
            \begin{scope}[xshift=1cm]
              \foreach \i in {1,...,3}
              \fill (\i*0.7, -2) circle (1pt);
            \end{scope}

            %% x y vecino.
            \begin{scope}[xshift=4cm]
              \node(3) [vertex] at (0, -2){};
              \node(4) [redV, label=270:$x$] at (1, -2){};
              \draw[edge] (3) to (4);
            \end{scope}

            %%  Nodos, esto se puede mejorar con un foreach.
            \begin{scope}[xshift=5cm]
              \node(5) [vertex] at (3, 1){};
              \node(6) [vertex] at (3, 0){};
              \foreach \i in {1,...,3}
              \fill (3, -\i*1) circle (1pt);
              \node(7) [vertex] at (3, -4){};
              \node(8) [vertex] at (3, -5){};
            \end{scope}

            %% Aristas de x hacia la derecha.
            \draw[edge] (4) to (5);
            \draw[edge] (4) to (6);
            \draw[edge] (4) to (7);
            \draw[edge] (4) to (8);

            %% Parte derecha.
            \begin{scope}[xshift=8cm]
              \foreach \j in {1,0,...,-5}
              \foreach \i in {1,...,3}
              \fill (\i*0.7, \j) circle (1pt);
              \node(9)  [vertex] at (3, 1){};
              \node(10) [blueV, label=0:$v_2$] at (4, 1){};
              \draw[edge] (9) to (10);
              \node(11) [vertex] at (3, 0){};
              \node(12) [blueV, label=0:$v_3$] at (4, 0){};
              \draw[edge] (11) to (12);
              \foreach \j in {3,4}
              \foreach \i in {-1,-2,-3}
              \fill (\j, \i) circle (1pt);
              \node(13)  [vertex] at (3, -4){};
              \node(14) [blueV, label=0:$v_{\Delta -1}$] at (4, -4){};
              \draw[edge] (13) to (14);
              \node(15) [vertex] at (3, -5){};
              \node(16) [blueV, label=0:$v_{\Delta}$] at (4, -5){};
              \draw[edge] (15) to (16);
            \end{scope}
          \end{tikzpicture}
        \end{figure}
      %%%%%%%%%%%%%%%%%%%%%%%%%%%%%%%%%%%%%%%%%%%%%%%%%%%%%%%%%%%%%%%%%%%%%%%%%%

        donde los $v_i$'s son las hojas, para $1 \leq i \leq \Delta$.
        \hfill $\square$
        \end{enumerate}

    %%%%%%%%%%%%%%%%%%%%%%%%%%%%%%% Ejercicio 3. %%%%%%%%%%%%%%%%%%%%%%%%%%%%%%%%
    \item Un {\em centro} en una gr\'afica es un v\'ertice
      $u$ tal que $\max_{v \in V} d(u, v)$ es m\'inima.
      Demuestre que un \'arbol tiene exactamente un centro
      o dos centros adyacentes.

      \renewcommand\qedsymbol{QED}
      \begin{proof}
        Sea $G$ un árbol.

        \textbf{Procedemos por contradicción.}

        Supongamos que $G$ tiene más de un centro tales que no son adyacentes.

        Sean $u, v \in V_{G}$ tales que $u$ y $v$ no son adyacentes.
        Esto significa que hay al menos un vértice entre ellos (llamémoslo $w$).

        Por definición de \textbf{centro}, sabemos que:
        $$max_{y \in V}d(x,y) \text{ es mínima}$$
        En nuestro caso, $d(u, v)$ es mínima.

        Pero notemos que $d(w, v) < d(u, v)$!!!, lo que implica que $w$ est\'a m\'as
        cerca del v\'ertice $v$ que el v\'ertice $u$.

        Como la contradicción yace de suponer que $u$ y $v$ no son adyacentes, podemos
        concluir que $u$ y $v$ son adyacentes.

        \textbf{Observación.} Veamos que $G$ no puede tener más de dos centros adyacentes.

        Supongamos que $G$ tiene tres centros adyacentes.

        Sean $w, v, x \in V_{G}$ tales que $w$ es adyacente a $x$ y $x$ es adyacente a $v$.

        Como vimos anteriormente, tendríamos que $d(x, v) < d(w, v)$ !!!.

        Esto contradice que el vértice $w$ sea un centro, ya que encontramos un vértice
        $x$ tal que su distancia a otro vértice $v$ es menor.
        La contradicción yace de suponer que $G$ tiene más de dos centros adyacentes.

        Por lo tanto, $G$ tiene a lo más dos centros adyacentes. \\
        \\
      \end{proof}

    %%%%%%%%%%%%%%%%%%%%%%%%%%%%%%% Ejercicio 4. %%%%%%%%%%%%%%%%%%%%%%%%%%%%%%%
    \item Demuestre o brinde un contraejemplo: Toda
      gr\'afica con menos aristas que v\'ertices tiene
      una componente que es un \'arbol.

      \renewcommand\qedsymbol{QED}
      \begin{proof}
        \textbf{Procedemos por contradicción.}

        Sea $G$ una gráfica tal que $|V_{G}| > |E_{G}|$.

        Supongamos que todas sus componentes conexas no son árboles.

        Esto es, $G_{1}, G_{2}, \dots, G_{n}$ componentes conexas de $G$ de donde
        cada $G_{i}$ con $i \in  \{1, 2, \dots, n\}$ no es un árbol.

        Como cada componente conexa de $G$ no es un árbol, entonces tienen
        ciclos.

        Notemos que por hipótesis, $|V_{G}| > |E_{G}|$. Lo que implica que existe
        al menos una componente conexa de $G$, digamos $G_{j}$ (con
        $j \in \{1, \dots, n\}$) tal que $G_{j} \cong K_{1}$.

        Por tanto, como $G_{j} \cong K_{1}$ y $K_{1}$ no tiene ciclos, la componente
        conexa $G_{j}$ es un árbol !!!.

        Lo anterior contradice el hecho de que toda componente conexa de $G$ no
        sea un árbol, pues si tienen ciclos significa que $|V_{G}| = |E_{G}|$.

        La contradicción yace de suponer que cada componente conexa de $G$ no es un
        árbol.

        Por lo tanto, podemos concluir que existe una componente conexa de $G$ que
        es un árbol. \\
        \\
      \end{proof}

    %%%%%%%%%%%%%%%%%%%%%%%%%%%%%%% Ejercicio 5. %%%%%%%%%%%%%%%%%%%%%%%%%%%%%%%
    \item Un {\em hidrocarburo saturado} es una
      mol\'ecula $C_mH_n$ en la que cada \'atomo de
      carbono tiene cuatro enlaces, cada
      \'atomo de hidr\'ogeno tiene un enlace, y
      ninguna sucesi\'on de enlaces forma un ciclo.
      Demuestre que para cualquier entero positivo
      $m$, la mol\'ecula $C_mH_n$ existe s\'olo si
      $n = 2m + 2$.

      \begin{proof}
        \textbf{Demostración por inducción sobre $m$.}

        \begin{itemize}
          \item \textbf{Paso base: $m = 1$}

            Por definición de hidrocarburo saturado $C_1 k_4 \Longrightarrow$  $4 = 2(1)+2$.

            Por lo tanto, para $m = 1$ se cumple que $n = 2m + 2$.

          \item \textbf{Hipótesis de Inducción: $m = k$}

            Si  $C_k H_n$, entonces supongamos $n = 2k + 2$.

          \item \textbf{Paso Inductivo:}

            Demostraremos para $m = k + 1$.

            Por \textbf{Hipótesis de Inducción}, tenemos que $C_k H_n \Longrightarrow
            n=2k+2$
            y por \textbf{Paso Base}, $C_1 K_4 \Longrightarrow 4=2(1)+2$.

            Sean $r$ que pertenece a los Naturales sin el $0$ y $C_r$, $C_1$ donde
            $C_r$ pertenece a $C_k H_n$ y $C_1$ pertenece a $C_1 k_4$ tal que $r$
            pertenece a $\{1, 2, 3, ..., k\}$.

            Si eliminemos $1$ hidrógeno a $C_r$ y $C_1 \Longrightarrow C_k H_{n-1}$
            y $C_1 k_3$ son iguales a:
            $$n-1 = 2k + 1 \text{ ...(1)}$$
            y
            $$3 = 2(1) + 1 \text{ ...(2)}$$
            Así, uniendo $C_k H_{n-1}$ y $C_1 k_3$ mediante los vértices $C_r$ y $C_1$,
            tenemos que $C_{k+1} H_{r}$ sería igual a la suma de (1) y (2). Esto es:
            $$n+2 = 2(k) + 2(1) + 2 \Longrightarrow  n+2 = 2(k+1) + 2 \Longrightarrow
            r = n + 2 \Longrightarrow  r = 2(k+1) + 2$$

            Por lo tanto, para $C_{k+1} H_{r}$ $r = 2(k+1) + 2$.

            Por lo anterior, podemos concluir que para todo $m$ que pertenece a
            Naturales sin el cero, $C_m H_n$ tal que $n = 2m + 2$.
        \end{itemize}
      \end{proof}

    %%%%%%%%%%%%%%%%%%%%%%%%%%%%%%% Ejercicio 6. %%%%%%%%%%%%%%%%%%%%%%%%%%%%%%%
    \item Demuestre que una sucesi\'on $(d_1, \dots,
      d_n)$ de enteros positivos es la sucesi\'on de
      grados de un \'arbol si y s\'olo si
      $\sum_{i=1}^n d_i = 2(n-1)$.

      \renewcommand\qedsymbol{QED}
      \begin{proof}
        Para este ejercicio veamos las dos implicaciones:
        \begin{itemize}
          \item[$\Rightarrow$)] Dada la sucesi\'on de grados $(d_1, \dots, d_n)$ de un
            \'arbol, entonces $\displaystyle \sum_{i=1}^n d_i = 2(n-1)$.

            Sabemos que para cualquier gr\'afica:
            \[
            \sum_{i=1}^n d_i = 2|E|
            \]
            Como en un \'arbol se cumple que $|E| = |V| - 1$, entonces $|E| = n - 1$. \footnote{Prop. $2.2.5$ en las
            notas de clase}

            Por tanto,
            \[
            \sum_{i=1}^n d_i = 2(n - 1)
            \]

          \item[$\Leftarrow$)] Dada una gr\'afica $G$ donde se cumple que
            $$\displaystyle \sum_{i=1}^n d_i = 2(n-1)$$
            , entonces $(d_1, \dots, d_n)$ es la sucesi\'on de grados en un \'arbol.

            Veamos que $G$ no tiene v\'ertices aislados.

            En el caso de tenerlos, supongamos sin p\'erdida de generalidad que
            $x \in V_G$ es un v\'ertice aislado. Entonces, $G - \{x\}$ no contiene
            v\'ertices aislados
            y $|E_{G - \{x\}}| = 2|V_{G - \{x\}}|$ !!

            Lo que implica que $G$ contiene como subgr\'afica inducida a alg\'un
            ciclo (pues la cantidad de v\'ertices ser\'ia de al menos la cantidad
            de v\'ertices) y se deja de cumplir que
            $$\displaystyle \sum_{i=1}^n d_i = 2(n-1)$$

            De la misma manera, podemos observar que todos los v\'ertices
            de $G$ no pueden tener al menos grado $2$, pues suponer esto nos genera
            al menos un ciclo e implicar\'ia que
            $$\displaystyle \sum_{i=1}^n d_i = 2(n) \text{!!!}$$
            , y claramente $2n \not= 2(n - 1)$.

            Luego, hay exactamente $2$ v\'ertices de grado $1$ en $G$, para esto
            llamemos $x, y \in V_G$ a los v\'ertices de grado $1$. Entonces, en
            $G -\{x, y\}$ se cumple que $\displaystyle \sum_{i=1}^n d_i = 2(n-2)$
            y al anexarle exactamente $\{x, y\}$ tendremos que
            \begin{eqnarray*}x
              \sum_{i=1}^n d_i &=& 2(n-2) + 2\\
              &=& 4n - 4 + 2\\
              &=& 4n - 2\\
              &=& 2(n - 1)
            \end{eqnarray*}

            En caso contrario, se deja de cumplir lo anterior y por esto
            se puede garantizar que estos son \'unicos.

            Como $G$ es ac\'iclica (por el argumento dado anteriormente), podemos
            deducir que $G$ es conexa. De no serlo, habr\'ia m\'as de $2$
            v\'ertices con grado $1$ y ya observamos que este no es un caso posible.

            Luego, $G$ es un \'arbol y por el ejercicio $1$ de esta tarea tenemos
            que $G$ es, part\'icularmente, una trayectoria.
        \end{itemize}
        Por lo tanto, podemos concluir que una sucesi\'on $(d_1, \dots,
        d_n)$ de enteros positivos es la sucesi\'on de grados de un \'arbol
        si y s\'olo si $\displaystyle \sum_{i=1}^n d_i = 2(n-1)$. \\
        \\
      \end{proof}
  \end{enumerate}

  %%%%%%%%%%%%%%%%%%%%%%%%%%%%%%% Extras %%%%%%%%%%%%%%%%%%%%%%%%%%%%%%%%%%%%%%%
  \section*{Puntos Extra}
  \begin{enumerate}
    %-------------------------------------- Extra 1.
    \item Para una gr\'afica conexa $G$ definimos la
      gr\'afica de \'arboles de $G$, $\mathcal{T}_G$,
      como la gr\'afica que tiene por v\'ertices a
      todos los \'arboles generadores de $G$, y tal que,
      si $S, T \in V_{\mathcal{T}_G}$, entonces $ST$ es
      una arista de $\mathcal{T}_G$ si y s\'olo si
      existen aristas $e \in E_S - E_T$ y $f \in E_T -
      E_S$ tales que $(S - e) + f = T$.   Demuestre que
      $\mathcal{T}_G$ es conexa.

      %-------------------------------------- Extra 2.
    \item Sea $T$ un \'arbol arbitrario con $k+1$ v\'ertices.
      Demuestre que si $G$ es simple y $\delta \ge k$,
      entonces $G$ tiene una subgr\'afica isomorfa a $T$.

      %-------------------------------------- Extra 3.
    \item Sea $\mathcal{T}$ una familia de sub\'arboles
      de un \'arbol $T$.   Deduzca, por inducci\'on sobre
      $|\mathcal{T}|$, que si cualesquiera dos elementos
      de $\mathcal{T}$ tienen un v\'ertice en com\'un,
      entonces hay un v\'ertice de $T$ que est\'a en
      todos los elementos de $\mathcal{T}$.

        %-------------------------------------- Extra 4.
    \item \begin{enumerate}
      \item Determine todos los arboles $T$ tales que
        $\overline T$ tambi\'en es un \'arbol.

        \textbf{\textit{Soluci\'on:}}
        \begin{itemize}
          \item[$\cdot$)]  Si $|V_T| = 1$, entonces por vacuidad se cumple
            el enunciado y terminamos.
          \item[$\cdot \cdot$)] Si $|V_T| > 1$, entonces sabemos existen $|V_T| -1$
            aristas para $T$ y que a lo más ${|V_T| \choose 2}$ si $T$ fuera
            completa.

            Notemos que $E_{\overline{T}} = {|V_T| \choose 2} - (|V_T| -1)$ y
            como queremos que $\overline{T}$ sea un árbol, entonces se debe cumplir
            la siguiente igualdad:
            \begin{eqnarray*}
              |V_T| -1 &=& {|V_T| \choose 2} - (|V_T| -1)\\
              \Leftrightarrow 2 \cdot (|V_T| -1) &=& {|V_T| \choose 2}\\
              \Leftrightarrow 2 \cdot (|V_T| -1) &=& \frac{|V_T| \cdot (|V_T| -1)}{2}\\
              \Leftrightarrow 4 \cdot \cancel{(|V_T| -1)} &=& |V_T| \cdot \cancel{(|V_T| -1)}\\
              \Rightarrow |V_T| &=& 4
            \end{eqnarray*}

            y del \textbf{Ejercicio $2$ de la Tarea $3$}, sabemos que hay $11$ gráficas
            de orden $4$ no isomorfas entre sí y sólo $2$ de esas son árboles.

            De \'estas dos \'ultimas, tenemos una es $P_4$ y la otra es el árbol
            tal que uno de sus vértices es de grado $3$. Pero en este último su
            complemento no es un árbol.

            Luego, $T = P_4$ y este es la único salvo isomorfismo.
            \hfill $\square$ \\
            \\
        \end{itemize}

      \item Determine todas las gr\'aficas de orden al
        menos cuatro tales que la subgr\'afica inducida por
        cualesquiera tres de sus v\'ertices es un \'arbol.

        \textbf{\textit{Soluci\'on:}}

        Veamos que si $|V| > 4$, se cumple cualquiera de las siguientes condiciones:
        \begin{itemize}
          \item[$\cdot$)] Si la gráfica es incompleta, existen al menos dos vértices
            que no se conectan mediante una arista y si no existe trayectoria entre estos,
            entonces existe una ``elección" de $3$ vértices tales que son inducidos de
            la gráfica original y no son un árbol.

          \item[$\cdot \cdot$)] Si la gráfica es completa, entonces hay $3$ vértices que
            al inducirlos en la gráfica original forman un $3$-ciclo y por tanto, no son un árbol.
        \end{itemize}
        De lo anterior la única gráfica que cumple con el enunciado es un $4$-ciclo
        que no tenga como gráfica inducida un $3$-ciclo.
        \hfill $\square$ \\
        \\
      \end{enumerate}
  \end{enumerate}
\end{document}
